\documentclass[a4paper]{article}

\usepackage{graphicx}
\usepackage{amsmath,amssymb}
\usepackage{minted}
\usepackage{multirow}
\usepackage{float}
\usepackage{todonotes}
\usepackage{forest}
\usepackage{xstring}
\usepackage{xcolor}
\usepackage[most]{tcolorbox}
\usepackage{xparse}

\usetikzlibrary{graphs,graphdrawing}
\usegdlibrary{layered}

% for external graphviz
\input{/home/donkarlo/Dropbox/repo/nd_ai_project/src/nd_ai/knoweldge_representation/language/formal/computer/programming/tex/latex/code_clips/external_graphviz.tex}

%for tags, comma separated \semtag{tag1, tag2}
\input{/home/donkarlo/Dropbox/repo/nd_ai_project/src/nd_ai/knoweldge_representation/language/formal/computer/programming/tex/latex/parts/control_sequence/kinds/semtag/semtag.tex}

% for links, just type \seehere
\input{/home/donkarlo/Dropbox/repo/nd_ai_project/src/nd_ai/knoweldge_representation/language/formal/computer/programming/tex/latex/code_clips/seehere.tex}

\title{}
\date{}

\begin{document}
    \maketitle
    According to Damasio's theory of consciousness, the protoself is the first stage in the hierarchical process of consciousness generation. Shared by many species, the protoself is the most basic representation of the organism, and it arises from the brain's constant interaction with the body. The protoself is an unconscious process that creates a "map" of the body's physiological state, which is then used by the brain to generate conscious experience. This "map" is constantly updated as the brain receives new stimuli from the body, and it forms the foundation for the development of more complex forms of consciousness.
    \url{https://en.wikipedia.org/wiki/Damasio%27s_theory_of_consciousness#Protoself}
    
    % \bibliographystyle{plain}
    % \bibliography{/home/donkarlo/Dropbox/repo/nd_shared_research/bibliography/bibliography.bib}
\end{document}